\subsection{HTTP Functionality}

The Initial prototype will send data via HTTP POST requests. The microcontroller is capable of creating these requests and sending them via the WiFi connection.This is a simple way of updating our web server with any new weight information. HTTP POST requests are simple and useful at doing what we need for this initial prototype.

\subsubsection{Request Format}

The request sends only a single line. The line is only the current weight measured by the sensor at the time of sending. Using a formatted string, the following is sent in the HTTP Packet: "weight=x.yz" Two points after the decimal are sent for accuracy for lower volume containers. 

The sensor data is converted to grams earlier in the code, so with an already calculated value all the HTTP code needs to do is package it and send it as a POST request. The post data is sent using the built in send\_http\_post() function. This function is a part of the microcontroller HTTP utilities. For production, HTTP data would be sent using C sockets. 

\subsubsection{Receiving Data}

On the server side, the POST data is received and using the expected format, decoded and sent to storage using Supabase. The data storage method and how the website pulls from the server will be discussed in their own section. 